% !PRINTED 386
\weberpage{386}

%<<CONT
so ist, da $-1:\omega'$ und $-1:\omega_1'$ auch reducirt sind, $\alpha$ die grösste in $-1:\omega_1'$ enthaltene ganze Zahl, und $\omega'$ ist nach (9) durch $\omega_1'$, also auch $\omega$ durch $\omega_1$ völlig bestimmt.

Hieraus folgt:

\gesperrt{3. Entwickelt man eine reducirte Zahl $\omega$ in einen Kettenbruch, so sind alle Schlusszahlen $\omega_n$ wieder reducirte Zahlen, und durch eine dieser Schlusszahlen $\omega_n$ ist sowohl die folgende $\omega_{n+1}$ als die vorangehende $\omega_{n-1}$ vollkommen bestimmt.}

\begin{center}
\S.\,125.

\gesperrt{Entwickelung reeller quadratischer Irrationalzahlen}\\
\gesperrt{in Kettenbrüche.}
\end{center}

\gesperrt{4. Für eine gegebene Discriminante giebt es nur eine endliche Anzahl reducirter Zahlen.}

Die Richtigkeit dieses Satzes sieht man leicht ein, wenn man nach \S.\,122 die Irrationalzahl $\omega$ in die Form setzt:

\weq{1}{\omega = \frac{b+\sqrt{D}}{2c} = \frac{-2a}{b-\sqrt{D}}, \quad D = b^2 + 4ac.}

Die conjugirte Zahl ist

\weq{2}{\omega' = \frac{b-\sqrt{D}}{2c} = \frac{-2a}{b+\sqrt{D}}.}

Soll $\omega$ reducirt sein, so ist

\weq{3}{0 < \frac{\sqrt{D}-b}{2c} < 1 < \frac{\sqrt{D}+b}{2c}.}

Es muss also $\sqrt{D}$ dasselbe Vorzeichen haben wie $c$; nehmen wir es positiv an, so ist

\weq{4}{0 < \sqrt{D} - b < 2c < \sqrt{D} + b.}

Eine zweite Form dieser Bedingungen erhalten wir aus der zweiten Darstellung

\weq{5}{0 < \sqrt{D} - b < 2a < \sqrt{D} + b.}

Es sind also $a$, $b$, $c$ positiv und $b$ kleiner als $\sqrt{D}$.

Bedeutet $\lambda$ die grösste in $\sqrt{D}$ enthaltene ganze Zahl, also
\[
\lambda < \sqrt{D} < \lambda + 1,
\]
%CONT>>
